% EM structure and teaching revision 20260928
\subsection{霍尔平衡与载流子浓度}
稳定载流子横向不再加速时，$q(\boldsymbol E_H+\boldsymbol v_d\times\boldsymbol B)=0$。电荷符号在这一步约去，但把漂移速度换为常规电流 $I=nqv_dwd$ 时，符号又通过 $q$ 进入测量关系。
取电流沿 $+x$、磁场沿 $+z$、宽度沿 $y$，定义 $U_H=V(-w/2)-V(w/2)$，有
\[
 E_y=v_dB,\qquad U_H=E_yw=\frac{IB}{nqd}.
\]
宽度消去，厚度 $d$ 留在分母。测出电压的符号可以判断单一载流子模型中的电荷符号，测出大小可以求浓度。
例如 $I=1\,\mathrm A$、$B=0.50\,\mathrm T$、$d=0.10\,\mathrm{mm}$，若 $|U_H|=1.0\,\mathrm{mV}$、$|q|=e$，则 $n\simeq3.1\times10^{25}\,\mathrm{m^{-3}}$。把毫米误作米会造成三个数量级的错误。
普通霍尔条形样品可在侧边积累电荷，圆环几何却没有同样的开路侧边；科尔比诺圆盘必须同时求径向和环向的运动。两种装置不能仅靠替换样品宽度互相套用公式。

